\documentclass[11pt,a4paper]{article}

\usepackage[T1]{fontenc}
\usepackage{lmodern}
\usepackage{amsmath,amssymb,amsthm}
\usepackage{mathrsfs}
\usepackage{graphicx}
\usepackage{hyperref}
\usepackage{microtype}
\usepackage{booktabs}
\usepackage{tabularx}
\usepackage{geometry}
\usepackage{doi}
\geometry{margin=2.5cm}

\input{glyphtounicode}
\pdfgentounicode=1

\hypersetup{
  pdftitle={Canonical Fourier Filtration and the Obstruction to a Spatial Continuum Limit},
  pdfauthor={Jérôme Beau},
  pdfsubject={The admissible fibre filtration of the Cosmochrony programme: exact identification and
    continuum-limit obstruction},
  pdfkeywords={Weil representation, Fourier filtration, continuum limit, Mosco convergence,
    no-go result, admissibility, non-injective projection}
}

\theoremstyle{plain}
\newtheorem{theorem}{Theorem}[section]
\newtheorem{proposition}[theorem]{Proposition}
\newtheorem{lemma}[theorem]{Lemma}
\newtheorem{corollary}[theorem]{Corollary}

\theoremstyle{definition}
\newtheorem{definition}[theorem]{Definition}
\newtheorem{hypothesis}[theorem]{Hypothesis}

\theoremstyle{remark}
\newtheorem{remark}[theorem]{Remark}

\newcommand{\Heis}{\mathrm{Heis}_3}
\newcommand{\Cq}{C_q}
\newcommand{\Gq}{G_q}
\newcommand{\T}{\mathbb{T}}
\newcommand{\Z}{\mathbb{Z}}
\newcommand{\R}{\mathbb{R}}
\newcommand{\Eq}{\mathcal{E}_q}
\newcommand{\CHeis}{C_{\mathrm{Heis}}}

\title{Canonical Fourier Filtration and the Obstruction to a Spatial Continuum Limit\\[0.5ex]
{\large What the Admissible Fibre Does and Does Not Converge To\\[0.3ex]
Q5a of the Cosmochrony Spectral Geometry Programme}}
\author{J\'er\^ome Beau\\[0.5ex]
\small Independent Researcher\\
\small \texttt{jerome.beau@cosmochrony.org}}
\date{2026}

\begin{document}

  \maketitle

  \begin{abstract}
    The companion Foundation paper establishes that the admissible fibre carries the Weil
    representation of $\Heis(\Z/q\Z)$ at a fixed non-trivial central character, and defines the
    canonical filtration of the fibre as the orbit spans
    $\Omega_n = \mathrm{span}\{\rho(g)v_0 : g \in B_n\}$ over breadth-first balls of the Cayley graph.
    This paper identifies that filtration exactly and determines what it does and does not converge to.

    Three results are established.
    First, an exact identification: with the pipeline's initial vector, $\Omega_n$ is precisely the
    toric Fourier window $\mathrm{span}\{e^{2\pi i b x/q} : |b| \le n\}$, of dimension $\min(2n+1, q)$;
    every fixed toric mode is captured once the published saturation depth $n_1(q)$ exceeds its index,
    while balanced (line-scale) profiles are rejected at all measured primes.
    Second, the published admissibility form converges to the zero form on this filtration, uniformly on
    norm-bounded sets, at rate $q^{-2}$.
    Third, a normalisation no-go: if the modulation and translation weights both remain positive and of
    order one --- as the pipeline data indicate --- then no common scalar normalisation of the form
    produces a non-trivial finite toric differential operator: preserving the derivative sector makes
    the modulation sector diverge, and preserving the modulation sector eliminates the derivative.
    The only non-trivial rescaled limit is a conditional Dirichlet operator in the rescaled frequency
    variable on the window itself, which does not provide a spatial continuum.
    Question Q5 of the Foundation programme therefore remains open, and the downstream identification of
    a flat spatial co-metric from a limit operator $-A\partial_x^2$ on $L^2(\R)$ rests on an input that
    is not established here.
    \textbf{Interpretive status.}
    The structural reading is that admissibility, as published, organises the fibre by frequency rather
    than by position: the emergent object is a growing Fourier window, not a discretised spatial line.
    Whether a spatial continuum can emerge from this structure is exactly the reopened content of Q5.
  \end{abstract}

  \paragraph{Keywords.}
    Weil representation; Fourier filtration; continuum limit; Mosco convergence; no-go result;
    admissibility; non-injective projection; toric windows

    \clearpage
    \tableofcontents

    \clearpage

  \section{Introduction}
  \label{sec:intro}

  \subsection{The question}

    The Foundation paper~\cite{Beau2026m} establishes, from four axioms governing admissible
    non-injective transitions, that the admissible fibre $F_n$ carries the Weil representation $V_\rho$
    of $\Gq = \Heis(\Z/q\Z)$ at a fixed non-trivial central character
    (Foundation Theorem; companion proof in~\cite{Beau2026HeisenbergStructure}), and defines, at its
    Level~3, the canonical filtration of the fibre by orbit spans over breadth-first (BFS) balls.
    Question Q5 of the programme asks how this discrete structure relates to continuous spacetime
    geometry in the large-$q$ limit.

    This paper answers a sharper preliminary question: \emph{what does the canonical filtration
    actually converge to, under the published admissibility form and its published normalisation?}
    The answer is negative in the direction Q5 hoped for, and exact.

  \subsection{Results and their statuses}

    \begin{description}
      \item[Exact identification (Theorem~\ref{thm:identification}; proved).]
      With the pipeline's initial vector $v_0$ (the uniform state), the canonical filtration stage is
      exactly the toric Fourier window
      $\Omega_n = \mathrm{span}\{e_b : |b| \le n\}$, $e_b(k) = e^{2\pi i b k/q}/\sqrt{q}$, of dimension
      $\min(2n+1, q)$.
      Every fixed toric mode $e_m$ lies in the stage as soon as the published saturation depth satisfies
      $n_1(q) \ge |m|$; balanced (line-scale) profiles are rejected at all measured primes
      (Section~\ref{sec:filtration}).

      \item[Zero-form theorem (Theorem~\ref{thm:zero}; proved).]
      With the published prefactor $q^{-2}$ and bounded weights, the admissibility form tends to zero
      uniformly on norm-bounded subsets of the filtration; its Mosco limit is the zero form.

      \item[Normalisation no-go (Theorem~\ref{thm:nogo}; proved, conditional on the measured weight
      behaviour).]
      If the modulation and translation weights have positive finite limits, then for every common
      scalar normalisation $\alpha_q$ the limit of the rescaled forms on fixed toric profiles is either
      the zero form, a bounded multiplication form with no derivative term, or the degenerate form
      $+\infty$ off the constants.
      No common scalar normalisation yields a non-trivial toric differential operator.

      \item[Dual window limit (Proposition~\ref{prop:blowup}; conditional).]
      Under the frequency blow-up $u = b/n_1(q)$, the rescaled form converges to a Dirichlet quadratic
      form on $L^2([-1,1])$ in the rescaled \emph{frequency} variable, with a potential term controlled
      by the critical coverage $x_1(q)$.
      This is a statement in the dual space; it does not produce a spatial continuum and does not
      resolve Q5.
    \end{description}

    Consequently, \emph{Q5 remains open}: the present analysis establishes what the admissible fibre
    converges to (a growing Fourier window; the zero form under the published normalisation; a
    dual-window Dirichlet form under blow-up) and what it does not converge to (any non-trivial toric
    differential operator under a common scalar normalisation, and in particular any operator of the
    form $-A\partial_x^2$ on a spatial $L^2$ space).
    The consequences for downstream papers are stated in Section~\ref{sec:consequences}.

  \section{Setup}
  \label{sec:setup}

  \subsection{Group, sector, and filtration}

    Fix a prime $q \ge 5$ and $\Gq = \Heis(\Z/q\Z)$ with generators $X, Y, Z$ (standard elementary
    matrices) and symmetric generating set $S_q = \{X^{\pm1}, Y^{\pm1}\}$; $B_n$ denotes the BFS ball
    of radius $n$ in the word metric.
    By the Bass--Guivarc'h theorem~\cite{Bass1972,Guivarch1973}, $|B_n| \sim \CHeis\, n^4$ with
    $\CHeis \approx 0.427$ for this generating set~\cite{Beau2026ccnote}.

    The representation space is $\Cq = L^2(\Z/q\Z)$ with counting norm $\|f\|^2 = \sum_k |f(k)|^2$.
    The Weil representation at central character $e^{2\pi i/q}$ (the corpus convention, $c = 1$) acts
    by
    \[
      (\rho(X) f)(k) = e^{2\pi i k/q} f(k), \qquad
      (\rho(Y) f)(k) = f(k+1 \bmod q), \qquad
      (\rho(Z) f)(k) = e^{2\pi i/q} f(k).
    \]
    The conjugate sector $c = q-1$ satisfies $\rho_{q-1} = \overline{\rho_1}$ and is identified with
    the $c = 1$ sector anti-linearly (the parity involution)~\cite{Beau2026a21,Beau2026a22}; it is
    never mixed linearly with it.

    The canonical filtration is the Foundation Level-3 orbit span with the pipeline's initial vector
    $v_0 = q^{-1/2}(1, \ldots, 1)$ (the uniform state, the convention of the O-series
    pipeline~\cite{Beau2026a16}):
    \begin{equation}
      \label{eq:filtration}
      \Omega_n \;=\; \mathrm{span}\{\rho(g)\, v_0 : g \in B_n\} \;\subseteq\; \Cq .
    \end{equation}

  \subsection{Published depth and weights, and their statuses}

    The published saturation depth $n_1(q)$ is the auto-calibrated marginal-novelty stop of the
    O-series exploration, with values $n_1 = 4, 8, 11, 13, 14$ at $q = 29, 61, 101, 151, 211$ and the
    exact reduction $n_1(q) = (x_1(q)/\CHeis)^{1/4}\sqrt{q}$ to the critical coverage
    $x_1(q) = |B_{n_1}|/q^2$~\cite{Beau2026a32,Beau2026ccnote}.

    \begin{hypothesis}[H-depth: unbounded depth]
      \label{hyp:depth}
      $n_1(q) \to \infty$ as $q \to \infty$.
    \end{hypothesis}

    Empirically $n_1$ increases monotonically ($4 \to 19$ over $q = 29 \to 601$) and
    $n_1(q)/q \to 0$ is established model-free~\cite{Beau2026a32}; Hypothesis~[H-depth] itself is
    unproved, and every capture statement below that needs it is conditional on it.
    Likewise, the empirical decay of the coverage ($x_1 \sim q^{-0.76}$ on the available range) is a
    numerical fit, not a demonstrated law; the asymptotic law of $x_1(q)$ is the open analytical
    programme of~\cite{Beau2026ccnote}.

    The admissibility weights $a_q(s)$, $s \in S_q$, are the fingerprint-response averages of the
    O-series pipeline~\cite{Beau2026a16,Beau2026a29}; the published admissibility form on the
    filtration stage is
    \begin{equation}
      \label{eq:form}
      \Eq(f) \;=\; \frac{q^{-2}}{2} \sum_{s \in S_q} a_q(s)\, \bigl\| (\rho(s) - I) f \bigr\|^2,
      \qquad f \in \Omega_{n_1(q)} ,
    \end{equation}
    with form domain the filtration stage (no interior projector).

    \begin{hypothesis}[H-$w'$: weight limits]
      \label{hyp:weights}
      $a_q(Y^{\pm1}) \to A_Y \in (0, \infty)$ and $a_q(X^{\pm1}) \to A_X \in (0, \infty)$.
    \end{hypothesis}

    Numerically, the four weight sequences lie in $[1.52, 2.01]$ across the measured primes with
    log-log slopes below $0.08$ in magnitude and exact $+/-$ symmetry; this supports~[H-$w'$] but does
    not prove it.
    All statements below that use~[H-$w'$] are conditional on it.

  \section{Exact identification of the filtration}
  \label{sec:filtration}

    \begin{theorem}[The canonical filtration is a Fourier window]
      \label{thm:identification}
      With $v_0$ the uniform state and $e_b(k) := e^{2\pi i b k/q}/\sqrt{q}$,
      \[
        \Omega_n \;=\; \mathrm{span}\bigl\{ e_b : |b| \le \min(n, \tfrac{q-1}{2}) \bigr\},
        \qquad
        \dim \Omega_n \;=\; \min(2n+1,\, q).
      \]
    \end{theorem}

    \begin{proof}
      For $g = (a, b, \gamma)$ in the coordinates of Section~\ref{sec:setup},
      $(\rho(g) v_0)(k) = e^{2\pi i(\gamma + b k)/q}\, v_0(k - a)
      = q^{-1/2}\, e^{2\pi i \gamma/q}\, e^{2\pi i b k/q}$:
      every orbit vector is the pure mode $e_b$ times a unimodular scalar.
      Hence $\Omega_n = \mathrm{span}\{e_b : b \in \beta(B_n)\}$ where $\beta(B_n)$ is the set of
      $b$-coordinates realised in $B_n$.
      Each generator changes the $b$-coordinate by at most one, so
      $\beta(B_n) \subseteq \{-n, \ldots, n\}$; conversely $Y^{b} \in B_n$ for $|b| \le n$, so equality
      holds.
      The modes $e_b$, $b$ ranging over distinct residues, are orthonormal, giving the dimension count.
    \end{proof}

    \begin{remark}[Physical bandwidth of the canonical stage]
      \label{rem:bandwidth}
      In balanced units ($h = \sqrt{2\pi/q}$), the stage $\Omega_{n_1(q)}$ has physical bandwidth
      $n_1(q)\,h = \sqrt{2\pi}\,(x_1(q)/\CHeis)^{1/4}$ exactly, by the reduction
      of~\cite{Beau2026ccnote}.
      The bandwidth is therefore governed by the critical coverage: it is bounded away from zero if and
      only if $x_1(q)$ is.
    \end{remark}

    \begin{corollary}[Mode capture and balanced rejection]
      \label{cor:capture}
      \textup{(i)} A fixed toric mode $e_m$ lies in $\Omega_{n_1(q)}$ if and only if
      $n_1(q) \ge |m|$; under Hypothesis~\textup{[H-depth]}, every fixed toric mode is captured
      permanently.
      \textup{(ii)} At the measured primes, balanced (line-scale) Hermite profiles are rejected: their
      relative distance to the stage is non-monotone in $q$ with a rising tail
      (for the Gaussian profile, $0.051,\, 0.010,\, 0.007,\, 0.010,\, 0.020$ over
      $q = 29, 61, 101, 151, 211$; for the sixth Hermite profile, flat at $0.75$--$0.74$), tracking the
      non-monotone physical bandwidth of Remark~\ref{rem:bandwidth}.
    \end{corollary}

    Part~(ii) is a reproducible numerical statement (script \texttt{code/canonical\_sector\_test.py},
    deterministic, published depths only); its status is empirical, at the measured primes.
    The structural mechanism is Remark~\ref{rem:bandwidth}: fixed toric modes require only depth, while
    balanced profiles require a bandwidth bounded away from zero, which the measured coverage decay
    contradicts.

  \section{The boundary of the window, component by component}
  \label{sec:boundary}

    On $\Omega_n$ the translation generator is diagonal, $\rho(Y) e_b = e^{2\pi i b/q} e_b$, and the
    modulation generator is the frequency shift, $\rho(X) e_b = e_{b+1}$.
    Three inequivalent operations must be distinguished at the window edges $b = \pm n$.

    \begin{lemma}[Boundary forms]
      \label{lem:boundary}
      Let $f = \sum_{|b| \le n} \hat f_b\, e_b \in \Omega_n$ with $2n + 1 < q$, and write
      $\Delta \hat f_b := \hat f_{b-1} - \hat f_b$.
      \begin{enumerate}
        \item[\textup{(a)}] \emph{Ambient restriction} (the convention of this paper, matching the form
        domain in~\eqref{eq:form}):
        \[
          \|(\rho(X) - I) f\|^2
          = \sum_{b=-n+1}^{n} |\Delta \hat f_b|^2 \;+\; |\hat f_{-n}|^2 \;+\; |\hat f_{n}|^2 ,
        \]
        the two edge terms arising from the outgoing shift ($\hat f_n$ feeding the external mode
        $e_{n+1}$) and the unmatched entry at $b = -n$.
        \item[\textup{(b)}] \emph{Operator compression}: the quadratic form of
        $P\,(\rho(X) - I)^*(\rho(X) - I)\,P = P\,(2I - \rho(X) - \rho(X)^{-1})\,P$, with $P$ the
        orthogonal projection onto $\Omega_n$, is identical to~\textup{(a)}: compression does not alter
        the diagonal entries, which remain $2$ at the edges.
        \item[\textup{(c)}] \emph{Edge deletion}: removing the outgoing edges (the graph Laplacian of
        the induced path, with degree $1$ at the endpoints) gives
        $\sum_{b=-n+1}^{n} |\Delta \hat f_b|^2$ with no edge terms.
        This modifies the operator, not merely the domain, and is not induced by any restriction or
        compression of the ambient operators.
      \end{enumerate}
      Under the frequency blow-up of Section~\ref{sec:blowup}, \textup{(a)} and \textup{(b)} impose
      Dirichlet boundary conditions at $u = \pm 1$, while \textup{(c)} would impose Neumann conditions.
      This paper adopts~\textup{(a)} throughout.
    \end{lemma}

    \begin{proof}
      (a) $(\rho(X) - I) f = \sum_{b=-n}^{n+1} (\hat f_{b-1} - \hat f_b)\, e_b$ with the conventions
      $\hat f_{\pm(n+1)} = 0$; expanding over the orthonormal modes gives the displayed sum, whose
      extreme terms are $|\hat f_{-n}|^2$ (at $b = -n$) and $|\hat f_n|^2$ (at $b = n+1$).
      (b) $P(2I - \rho(X) - \rho(X)^{-1})P$ has matrix entries $2$ on the diagonal of the whole window
      and $-1$ between neighbours inside the window; its quadratic form telescopes to exactly the same
      expression as~(a).
      (c) The path Laplacian has diagonal degree $1$ at the endpoints; its form is the interior sum
      only.
      For the boundary conditions: with $\hat f_b = n^{-1/2}\varphi(b/n)$ and the blow-up prefactor
      $n^{2}$ of Section~\ref{sec:blowup}, the edge terms scale as $n\,|\varphi(\pm 1)|^2 \to \infty$
      unless $\varphi(\pm 1) = 0$ (Dirichlet), while their absence in~(c) leaves the natural (Neumann)
      condition.
    \end{proof}

  \section{The published normalisation: the zero form}
  \label{sec:zero}

    \begin{theorem}[Zero-form limit]
      \label{thm:zero}
      Assume only $\sup_q \max_{s \in S_q} a_q(s) < \infty$.
      Then the published form~\eqref{eq:form} satisfies
      \[
        \Eq(f) \;\le\; C\, q^{-2}\, \|f\|^2 \qquad \text{for all } f \in \Omega_{n_1(q)},
      \]
      with $C = 8 \sup_q \max_s a_q(s)$.
      Consequently $\Eq \to 0$ uniformly on norm-bounded sets, and the Mosco limit of the published
      forms, along any isometric identification of the stages with subspaces of a fixed Hilbert space,
      is the zero form.
    \end{theorem}

    \begin{proof}
      $\|(\rho(s) - I) f\| \le 2\|f\|$ for unitary $\rho(s)$; there are four generators; the prefactor
      is $q^{-2}/2$.
      The liminf inequality for the zero form is trivial, and any bounded recovery sequence gives
      $\limsup \Eq(f_q) = 0$.
    \end{proof}

    The zero-form conclusion is thus not a marginal effect of the filtration: it follows from the
    published prefactor alone, on the whole space as on the stages.

  \section{The normalisation no-go}
  \label{sec:nogo}

    Write the bare weighted form
    $\mathcal{B}_q(f) = \sum_{s \in S_q} a_q(s)\, \|(\rho(s) - I) f\|^2$, so that
    $\Eq = (q^{-2}/2)\,\mathcal{B}_q$, and consider a general common scalar normalisation
    $\alpha_q\,\mathcal{B}_q$, $\alpha_q > 0$.
    Fix a smooth toric target $f$ with finitely many Fourier modes and let
    $f_q \in \Omega_{n_1(q)}$ be its mode truncation (exact once $n_1(q)$ exceeds the largest mode
    index, hence eventually exact under~[H-depth]).
    Then, exactly,
    \[
      \mathcal{B}_q(f_q)
      = \bigl(a_q(X) + a_q(X^{-1})\bigr)\, c_X(f)
      + \bigl(a_q(Y) + a_q(Y^{-1})\bigr) \sum_b 4 \sin^2(\pi b/q)\, |\hat f_b|^2 ,
    \]
    with $c_X(f) = \sum_b |\Delta \hat f_b|^2 = \int_\T |e^{2\pi i x} - 1|^2\, |f(x)|^2\,\mathrm{d}x
    > 0$ for non-constant $f$, and the $Y$-sum equal to
    $q^{-2}\bigl[(2\pi)^2 \sum_b b^2 |\hat f_b|^2\bigr](1 + o(1))$.

    \begin{theorem}[No common scalar normalisation yields a toric differential limit]
      \label{thm:nogo}
      Assume Hypotheses~\textup{[H-depth]} and~\textup{[H-$w'$]}.
      Let $\alpha_q > 0$ be any sequence of common scalar normalisations.
      Then every subsequential limit of the rescaled forms $\alpha_q\,\mathcal{B}_q$ on smooth toric
      targets is one of:
      \begin{enumerate}
        \item[\textup{(i)}] the zero form (when $\alpha_q \to 0$);
        \item[\textup{(ii)}] the bounded multiplication form
        $2\alpha_\ast A_X \int_\T 2\bigl(1 - \cos 2\pi x\bigr)\, |f|^2$ (when
        $\alpha_q \to \alpha_\ast \in (0, \infty)$), which contains no derivative term;
        \item[\textup{(iii)}] the degenerate form equal to $+\infty$ on every non-constant target (when
        $\alpha_q \to \infty$).
      \end{enumerate}
      In particular, no choice of $\alpha_q$ produces a limit containing both a finite derivative term
      and the modulation sector: preserving the derivative sector ($\alpha_q \sim q^2$) makes the
      modulation sector diverge, and preserving the modulation sector ($\alpha_q \to \alpha_\ast$)
      eliminates the derivative.
    \end{theorem}

    \begin{proof}
      Every sequence $\alpha_q$ has subsequences converging in $[0, \infty]$; on each, the displayed
      exact decomposition gives:
      $\alpha_q\,\mathcal{B}_q(f_q) \to 0$ if $\alpha_q \to 0$ (both sectors bounded);
      $\to 2\alpha_\ast A_X\, c_X(f)$ if $\alpha_q \to \alpha_\ast$ (the $Y$-sector carries the factor
      $q^{-2}$);
      $\to +\infty$ if $\alpha_q \to \infty$ and $c_X(f) \ne 0$, i.e.\ for every non-constant $f$.
      The multiplication identification in~(ii) is Parseval:
      $c_X(f) = \int_\T |e^{2\pi i x} - 1|^2\, |f|^2 = \int_\T 2(1 - \cos 2\pi x)\, |f|^2$.
    \end{proof}

    \begin{remark}[The multiplication limit is not geometric]
      \label{rem:multiplication}
      The regime~(ii) is the only non-zero limit reachable with the derived weights and a bounded
      normalisation.
      Its limit operator is multiplication by $4 \alpha_\ast A_X (1 - \cos 2\pi x)$: a bounded
      operator, with no differential symbol of any positive order, hence no principal symbol, no
      co-metric, and no diffusion.
      It cannot serve as the generator of an emergent spatial geometry.
    \end{remark}

    \begin{remark}[Anisotropic renormalisation is not corpus-derived]
      \label{rem:aniso}
      A per-generator renormalisation (e.g.\ $q^2 a_q(Y^{\pm1})$ against $a_q(X^{\pm1})$) would retain
      both sectors.
      No such structure is derived anywhere in the corpus: the weights form a single family defined
      symmetrically over the generating set, the published prefactor is common, and the BFS word metric
      treats the generators symmetrically.
      Introducing it would be a new input, not a derivation, and is not done here.
    \end{remark}

  \section{The conditional dual window limit}
  \label{sec:blowup}

    The filtration's own natural profiles are envelopes over the window: for
    $\varphi \in C^\infty([-1,1])$ set
    $f_q = \sum_{|b| \le n_1} n_1^{-1/2}\, \varphi(b/n_1)\, e_b$ with $n_1 = n_1(q)$.

    \begin{proposition}[Dual blow-up limit]
      \label{prop:blowup}
      Assume~\textup{[H-depth]} and~\textup{[H-$w'$]}, and adopt the ambient-restriction boundary
      convention of Lemma~\ref{lem:boundary}\textup{(a)}.
      With the common prefactor $\alpha_q = n_1(q)^2$, and for envelopes with
      $\varphi(\pm 1) = 0$,
      \[
        \alpha_q\, \mathcal{B}_q(f_q)
        \;\longrightarrow\;
        2 A_X \int_{-1}^{1} |\varphi'(u)|^2\, \mathrm{d}u
        \;+\;
        2 A_Y\, (2\pi)^2\, \lambda \int_{-1}^{1} u^2\, |\varphi(u)|^2\, \mathrm{d}u ,
      \]
      along any subsequence on which $x_1(q)/\CHeis \to \lambda \in [0, \infty)$, using the exact
      identity $(n_1^2/q)^2 = x_1/\CHeis\,(1 + o(1))$; envelopes with $\varphi(\pm 1) \neq 0$ are
      excluded in the limit (Dirichlet), by Lemma~\ref{lem:boundary}.
      Under the empirical coverage decay $\lambda = 0$ and the limit is the Dirichlet form of
      $-2A_X\,\partial_u^2$ on $H^1_0([-1,1])$; if $x_1(q)$ were bounded away from zero, the limit
      would be an oscillator-type Schr\"odinger form in the dual variable.
    \end{proposition}

    \begin{proof}[Proof sketch]
      The $X$-sector is the discrete Dirichlet energy of the envelope on the window
      (Lemma~\ref{lem:boundary}(a)); with spacing $1/n_1$ and the prefactor $n_1^2$ it converges to
      the Dirichlet integral, the edge terms enforcing $\varphi(\pm 1) = 0$.
      The $Y$-sector is diagonal with symbol $4\sin^2(\pi b/q) = (2\pi b/q)^2(1 + o(1))$ on the
      window; on envelopes it equals $(2\pi n_1/q)^2 \int u^2 |\varphi|^2\, (1 + o(1))$, and
      $n_1^2 \cdot (n_1/q)^2 = (n_1^2/q)^2$.
    \end{proof}

    \begin{remark}[What this limit is, and is not]
      \label{rem:dual-status}
      Proposition~\ref{prop:blowup} is a statement in the rescaled \emph{frequency} variable
      $u = b/n_1(q)$, on the window $[-1,1]$.
      It is the unique non-trivial limit reachable with the derived weights and a single common
      prefactor, and its constants are exact; but its underlying variable is dual, not spatial.
      It therefore does not resolve Q5 and must not be read as substituting the frequency variable for
      the spatial one.
      Its conditional status is threefold: [H-depth], [H-$w'$], and --- for the presence or absence of
      the potential term --- the open asymptotic law of $x_1(q)$.
    \end{remark}

  \section{Consequences}
  \label{sec:consequences}

  \subsection{Q5 remains open}

    The programme's Question Q5 asked for a derivation of continuous spatial structure from the
    admissible fibre.
    The present results replace that expectation with an exact statement:
    the canonical filtration is a growing toric Fourier window; the published form converges to zero on
    it; no common scalar normalisation produces a non-trivial toric differential operator; and the only
    non-trivial rescaled limit lives in the dual variable.
    \emph{Q5 is therefore open}: no spatial continuum limit is established by the admissibility
    structure as published.

  \subsection{Downstream dependency}

    The companion paper Q5b~\cite{Beau2026q5b} consumes a limit operator $-A\partial_x^2$ on $L^2(\R)$
    and reads an effective co-metric from its principal symbol.
    That input is not established by the present analysis: none of the derived limits (zero form,
    multiplication form, dual-window Dirichlet form) is a second-order operator in a spatial variable.
    Q5b's construction is accordingly conditional on an input whose derivation is open.

  \subsection{A conditional route kept on record}

    A balanced-scale analysis (sample spacing $q^{-1/2}$), under which both generators rescale to the
    canonical pair $(x, -i\partial_x)$ and the form converges to an oscillator-type operator with a
    spectral gap, remains available as a \emph{conditional possibility}: by
    Remark~\ref{rem:bandwidth} it requires the physical bandwidth of the canonical stage to stay
    bounded away from zero, i.e.\ a future non-vanishing asymptotic bound on the critical coverage
    $x_1(q)$.
    The measured coverage decay currently contradicts this; the question coincides with the open
    programme of~\cite{Beau2026ccnote} and is recorded as such, not as a result.

  \subsection{Status table}

  \begin{table}[h]
    \centering
    \small
    \renewcommand{\arraystretch}{1.3}
    \begin{tabularx}{\textwidth}{@{}lXp{3.6cm}@{}}
        \toprule
        Result & Content & Status \\
        \midrule
        Thm~\ref{thm:identification} & $\Omega_n$ = Fourier window, $\dim = \min(2n{+}1, q)$
        & proved \\
        Cor~\ref{cor:capture}(i) & permanent capture of fixed modes & conditional on [H-depth] \\
        Cor~\ref{cor:capture}(ii) & balanced-profile rejection & numerical, measured primes \\
        Lem~\ref{lem:boundary} & ambient restriction and compression give Dirichlet; edge deletion
        gives Neumann & proved \\
        Thm~\ref{thm:zero} & published form $\to 0$, rate $q^{-2}$ & proved \\
        Thm~\ref{thm:nogo} & no common scalar normalisation gives a toric differential limit
        & proved under [H-depth], [H-$w'$] \\
        Rem~\ref{rem:aniso} & anisotropic renormalisation would be a new input & corpus audit \\
        Prop~\ref{prop:blowup} & dual window Dirichlet limit, exact constants
        & conditional ([H-depth], [H-$w'$], $x_1$ law) \\
        Q5 & spatial continuum from admissibility & open \\
        \bottomrule
    \end{tabularx}
    \caption{Status of all results.
    [H-depth] and [H-$w'$] are numerically supported hypotheses (Section~\ref{sec:setup}); no
    numerical statement is promoted to a theorem.}
    \label{tab:status}
  \end{table}

  \subsection{Open problems}

  \noindent\textbf{Q5a-O1} (Coverage law).
  Derive the asymptotic law of the critical coverage $x_1(q)$ --- equivalently of the physical
  bandwidth of the canonical stage.
  This is the single quantity on which both the dual potential term and any revival of a balanced
  spatial picture depend; it coincides with the open programme of~\cite{Beau2026ccnote}.

  \medskip
  \noindent\textbf{Q5a-O2} (Depth divergence).
  Prove Hypothesis~[H-depth]: $n_1(q) \to \infty$.

  \medskip
  \noindent\textbf{Q5a-O3} (Weights).
  Prove Hypothesis~[H-$w'$] analytically from the O-series spectral bounds, replacing its current
  numerical support.

  \medskip
  \noindent\textbf{Q5a-O4} (Boundary selection).
  Determine whether the corpus axioms select the ambient-restriction convention of
  Lemma~\ref{lem:boundary}(a) (Dirichlet) or another boundary operation; the dual limit's boundary
  condition depends on it.

  \medskip
  \noindent\textbf{Q5a-O5} (Meaning of the dual window).
  Interpret the dual-variable Dirichlet limit of Proposition~\ref{prop:blowup} within the programme,
  or show that it is a structural artefact of the filtration; in either case without substituting it
  for the missing spatial limit.

  \section{Conclusion}
  \label{sec:conclusion}

    The canonical filtration of the admissible fibre is exactly a growing toric Fourier window.
    With the published form and normalisation, the limit is the zero form; if the modulation and
    translation weights both remain positive and of order one, no common scalar normalisation produces
    a finite non-trivial toric differential operator --- preserving the derivative sector makes the
    modulation sector diverge, while preserving the modulation sector eliminates the derivative.
    The only non-trivial rescaled limit is a conditional Dirichlet operator on the rescaled frequency
    window, a dual-space statement.
    Q5 therefore remains open, and downstream constructions that consume a spatial operator
    $-A\partial_x^2$ rest on an input that is not established.

    \medskip
    \textbf{Interpretive outlook.}
    The structural reading of these results is that the admissibility structure, as published,
    organises the fibre by \emph{frequency}: what grows with the exploration is a bandwidth, not a
    spatial resolution.
    A spatial continuum would require either a demonstrated non-vanishing asymptotic bandwidth (the
    open coverage law) or a structure the corpus has not yet derived (an anisotropic normalisation, or
    a different filtration).
    Making that requirement precise --- rather than assuming it --- is what the present negative
    result contributes to Q5.

  \appendix

  \section{Weil representation: explicit generator actions}
  \label{app:weil}

  For $f \in \Cq = L^2(\Z/q\Z)$ and $k \in \Z/q\Z$, at central character $e^{2\pi i/q}$:
  \begin{align*}
    (\rho(X) f)(k) &= e^{2\pi i k/q} f(k) \qquad \text{(modulation; frequency shift on modes)},\\
    (\rho(Y) f)(k) &= f(k+1 \bmod q) \qquad \text{(cyclic shift; diagonal on modes)},\\
    (\rho(Z) f)(k) &= e^{2\pi i/q} f(k) \qquad \text{(central character)}.
  \end{align*}
  The discrete Stone--von Neumann theorem guarantees uniqueness (up to isomorphism) of the irreducible
  representation at this character, underpinning the identification $F_n \simeq V_\rho$ of
  Foundation~\cite{Beau2026m,Beau2026HeisenbergStructure}.

  \bibliographystyle{plain}
  \bibliography{cosmochrony-bibliography,external-refs}

\end{document}
